% !TEX root = ../../main.tex

\section{Resultados generales}
\label{sec:resultados-generales}

La metodología de validación de la solución mostrado en la Figura~\ref{fig:validacion-flujo} se aplicó a 93 programas que no correspondían a controles de activación: 44 del corpus de \texttt{Maybe} y 49 del corpus probabilístico. En todos ellos, la variante de \texttt{ApplicativeDo}, la selección automática y cada selección forzada conservaron tanto el código de salida como la salida observable del programa original. En particular, los casos de \texttt{Dist.T Rational} produjeron exactamente la misma distribución después de aplicar \texttt{Dist.norm}, incluidos aquellos con pesos no uniformes, resultados repetidos, soportes dependientes y condicionamiento. Por tanto, todos los casos evaluados respondieron favorablemente \textbf{PV1} dentro del criterio observacional establecido.

En los mismos 93 programas, el costo de cada candidato compilado individualmente permitió recalcular el mínimo, su primer índice y la cantidad de empates. En todos los casos, estos valores coincidieron con las métricas emitidas por el Renamer y el candidato elegido automáticamente correspondió al primer plan de costo mínimo entre los reordenamientos válidos generados. Este resultado responde favorablemente \textbf{PV2}. Además, como el orden fuente siempre pertenece al conjunto de candidatos, la selección automática nunca produjo un costo mayor que el obtenido por \texttt{ApplicativeDo} sobre ese mismo orden.

Los seis programas restantes asociados al control de activación de la extensión respondieron \textbf{PV3} y también produjeron el comportamiento esperado. En ambos corpus, los cuatro programas con bloques \texttt{do} ordinarios mantuvieron el marcador no activado, mostrado en el Renamer como \texttt{commutative-do = False}, incluso cuando \texttt{QualifiedDo} estaba habilitado, la interfaz conmutativa estaba importada y se había declarado la instancia correspondiente. Los dos bloques \texttt{do} escritos mediante \texttt{CD.do} produjeron \texttt{commutative-do = True} a travez de las trazas del Renamer, cada uno generando seis candidatos, incluyendo el orden original. En consecuencia, el uso completo de la solución requiere \texttt{ApplicativeDo}, \texttt{QualifiedDo}, la interfaz calificada, una instancia \texttt{CommutativeMonad} y el marcador local \texttt{CD.do} en el bloque que quiera ser declarado como conmutativo. La instancia expresa el contrato de conmutatividad y se comprueba durante el tipado, mientras que el marcador resoluble es la condición que distingue directamente ambas rutas dentro del Renamer.

La cantidad de casos que reducen estrictamente el costo de \texttt{ApplicativeDo} no se utiliza como indicador agregado de efectividad. La mayor parte del corpus fue construida para comprobar preservación observable, formación de dependencias, recolección de binders o selección correcta, y no para disponer deliberadamente un orden fuente desfavorable. Por ello, es esperable que una cadena de statements admita un único orden, que varios reordenamientos válidos tengan el mismo costo o que \texttt{ApplicativeDo} ya alcance el costo mínimo. Los cinco programas de cada corpus que obtuvieron mejores resultados respecto a \texttt{ApplicativeDo} fueron tres casos de la familia \texttt{cost-selection}, el caso asociado a dos cadenas de statements independientes y el caso en el que se usa \texttt{let} independiente.

Establecidos estos resultados globales, las secciones siguientes presentan casos representativos de las familias evaluadas. El objetivo es mostrar cómo los fenómenos presentes en el código fuente se reflejan en el grafo de precedencia RAW y cómo esa estructura determina los reordenamientos y planes que puede considerar la extensión.
